\documentclass[a4paper,11pt]{article}
\usepackage[utf8]{inputenc}
\usepackage[T1]{fontenc}
\usepackage[english]{babel}
\usepackage[left=1cm,right=1cm,top=.5cm,bottom=0cm]{geometry}
\usepackage{enumitem}
\usepackage{hyperref}
\usepackage{xcolor}
\usepackage{titlesec}
\usepackage{marvosym}
\usepackage{lmodern}
\usepackage[normalem]{ulem}

\definecolor{primary}{RGB}{0, 129, 145}
\definecolor{secondary}{RGB}{80, 80, 80}
\hypersetup{colorlinks=true, linkcolor=primary, filecolor=primary, urlcolor=primary}
\titleformat{\section}{\large\bfseries\color{primary}\uppercase}{}{0em}{}[\titlerule]
\titlespacing{\section}{0pt}{8pt}{5pt}

\begin{document}
\begin{center}
    {\Large \textbf{Lo\"ic Aron MBASSI EWOLO}} \\ \vspace{4pt}
    \textbf{\large Alternant Ing\'enieur IoT -- Syst\`emes Connect\'es \& Smart Metering} \\
    \textit{\small Stage INRIA Saclay (Avril--Ao\^ut 2026) -- Disponible en alternance d\`es septembre 2026 (12 mois)} \\ \vspace{4pt}
    \small
    \Pointinghand\ Cergy, France (Mobile IDF) \quad \Mobilefone\ (+33) 7 59 52 33 98 \quad
    {\color{primary}\Letter\ \href{mailto:loicaron.mbassiewolo@ensea.fr}{loicaron.mbassiewolo@ensea.fr}} \\ \vspace{2pt}
    {\color{primary}LinkedIn: \href{https://www.linkedin.com/in/loic-mbassi/}{linkedin.com/in/loic-mbassi}} \quad
    {\color{primary}GitHub: \href{https://github.com/Nameless0l}{github.com/Nameless0l}} \quad
    {\color{primary}Portfolio: \href{https://mbassiloic.tech}{mbassiloic.tech}}
\end{center}
\vspace{-6pt}
\noindent\small{Conception et d\'eploiement de syst\`emes IoT complets (capteurs, r\'eseaux LPWAN, backend), avec maîtrise des protocoles smart metering et optimisation d'inf\'erence ML embarqu\'ee. Formation double dipl\^ome ENSEA/Polytechnique Yaound\'e, sp\'ecialit\'e syst\`emes embarqu\'es et IA. Stage INRIA Saclay (PRIMO) en cours sur pipeline Python avanc\'e et analyse statique.}

\section{Formation}
\begin{itemize}[leftmargin=*, label=\tiny$\bullet$, nosep]
    \item \textbf{ENSEA -- \'Ecole Nationale Sup\'erieure de l'\'Electronique et de ses Applications} \hfill \textit{2025 -- 2027} \\
    \small{Double dipl\^ome d'Ing\'enieur -- Syst\`emes Embarqu\'es \& Intelligence Artificielle.}
    \item \textbf{\'Ecole Polytechnique de Yaound\'e (1\textsuperscript{\`ere} \'ecole d'ing\'enieur CEMAC)} \hfill \textit{2021 -- 2025} \\
    \small{Dipl\^ome d'Ing\'enieur de Conception (Master 2) : \textbf{Math\'ematiques Appliqu\'ees}, Algorithmique, G\'enie Logiciel, mention Tr\`es Bien.}
\end{itemize}

\section{Comp\'etences Techniques}
\begin{itemize}[leftmargin=*, nosep]
    \item \small{\textbf{IoT/R\'eseaux :} LPWAN (LoRa, Sigfox), Protocoles IoT, Smart Metering, Architecture Syst\`emes Connect\'es, Acquisition/Traitement Capteurs, \'Energie Bas-d\'ebit.}
    \item \small{\textbf{Embarqu\'e/Hardware :} STM32, C, I2C, SPI, UART, Capteurs (Temp, Humidit\'e, Pr\'esence), TensorFlow Lite, Oscilloscope, Soudure, Prototypage.}
    \item \small{\textbf{ML/Data Science :} Python, TensorFlow, scikit-learn, pandas, numpy, D\'etection Anomalies, S\'eries Temporelles, Pipelines ETL, \'Evaluation Mod\`eles.}
    \item \small{\textbf{Syst\`emes :} VHDL, FPGA, Linux Embarqu\'e, Docker, Git, Tests Unitaires, Architecture Microservices, MQTT.}
    \item \small{\textbf{Langues :} Fran\c{c}ais natif, Anglais courant professionnel, Notions allemand.}
\end{itemize}

\section{Exp\'eriences \& Projets}
\begin{itemize}[leftmargin=*, label={-}]
    \item \textbf{Stage INRIA Saclay -- \'Equipe TRIBE} \hfill \textit{Avril--Ao\^ut 2026} \\
    \textit{Palaiseau (IP Paris)} \hfill \textit{Python, NLP, BERT, Transformers, scikit-learn, Analyse Statique}
    \vspace{-2pt}
        \begin{itemize}[leftmargin=0.4cm, nosep, label=\tiny$\bullet$]
            \item \small{Projet PRIMO : pipeline NLP pour analyse confiance/transparence apps mobiles (Louvre, SNCF), RGPD compliance.}
            \item \small{Analyse statique SDKs embarqu\'es, identification risques fuite donn\'ees localisation et personnelles.}
            \item \small{Scoring interpr\'etable (scikit-learn), \'evaluation mod\`eles BERT/Transformers \'a grande \'echelle.}
            \item \small{Collaboration \'Ecole des Ponts, programme PEPR MOBIDEC, architecture logicielle avanc\'ee.}
        \end{itemize}
    \vspace{3pt}
    \item \textbf{Stage IoT Lab ENSPY -- Optimisation ML pour Hardware Embarqu\'e} \hfill \textit{Oct 2023--Jan 2024} \\
    \textit{Yaound\'e, Cameroun} \hfill \textit{Python, TensorFlow Lite, Arduino, Raspberry Pi, Capteurs, C}
    \vspace{-2pt}
        \begin{itemize}[leftmargin=0.4cm, nosep, label=\tiny$\bullet$]
            \item \small{Optimisation mod\`eles ML pour contraintes hardware limit\'e, quantisation, pruning, compression aggressive.}
            \item \small{D\'eploiement TensorFlow Lite sur microcontr\^oleurs (Arduino, Raspberry Pi Zero), inférence temps r\'eel.}
            \item \small{Acquisition donn\'ees capteurs, pipelines temps r\'eel, optimisation latence et consommation \'energie.}
            \item \small{Validation oscilloscope, calibration capteurs, documentation syst\`eme complet.}
        \end{itemize}
    \vspace{3pt}
    \item \textbf{Harmony Gloves -- Capteurs IMU/Flex \& Acquisition I2C/SPI sur STM32} \hfill \textit{2024} \\
    \textit{Labo IOTA ENSEA, Projet \'Equipe} \hfill \textit{STM32, I2C, SPI, Capteurs, Oscilloscope, Soudure PCB}
    \vspace{-2pt}
        \begin{itemize}[leftmargin=0.4cm, nosep, label=\tiny$\bullet$]
            \item \small{Participation conception PCB custom, soudure capteurs flex et IMU (MPU6050), test d\'everminage oscilloscope.}
            \item \small{Acquisition donn\'ees via I2C/SPI sur STM32, calibration capteurs, traitement signal temps r\'eel.}
            \item \small{D\'eploiement TinyML pour classification gestuelle, optimisation puissance embarqu\'ee (basse consommation).}
            \item \small{Test robustesse syst\`eme multiples sc\'enarios, documentation mat\'erielle et logicielle.}
        \end{itemize}
    \vspace{3pt}
    \item \textbf{Urban Waste Management -- Architecture IoT Compl\`ete \& ML Pr\'edictif} \hfill \textit{2024} \\
    \textit{1\textsuperscript{er} Prix CITS National (vs 75 \'equipes)} \hfill \textit{Python, ML, IoT, Capteurs, LPWAN, Backend}
    \vspace{-2pt}
        \begin{itemize}[leftmargin=0.4cm, nosep, label=\tiny$\bullet$]
            \item \small{Architecture IoT : capteurs capacitifs pour mesure remplissage, r\'eseau LPWAN, backend centralis\'e.}
            \item \small{Mod\`ele ML (scikit-learn, XGBoost) pr\'ediction remplissage conteneurs, d\'etection anomalies temps r\'eel.}
            \item \small{Optimisation collecte d\'echets (algo TSP), r\'eduction -20\% co\^uts transport et \'emissions CO2.}
            \item \small{Prototype fonctionnel d\'eploy\'e, dashboards temps r\'eel, 1\textsuperscript{er} prix national vs 75 \'equipes.}
        \end{itemize}
    \vspace{3pt}
    \item \textbf{Conception FPGA -- Syst\`eme SoPC Cyclone V \& Nios V} \hfill \textit{2025--en cours} \\
    \textit{ENSEA} \hfill \textit{VHDL, Quartus, Nios V, I2C, JTAG, Debug, Architecture RTL}
    \vspace{-2pt}
        \begin{itemize}[leftmargin=0.4cm, nosep, label=\tiny$\bullet$]
            \item \small{Description RTL en VHDL, simulation comportementale, intégration contraintes routage (Quartus).}
            \item \small{Soft-processeur Nios V, contr\^oleur I2C pour capteur ADXL345, debug JTAG complet.}
            \item \small{Validation oscilloscope, optimisation timing, synth\`ese syst\`eme complet sur Cyclone V.}
            \item \small{Interfaces protocoles standards (I2C, SPI), architecture scalable pour extensions futures.}
        \end{itemize}
    \vspace{3pt}
    \item \textbf{Assistant de Recherche -- MILA, \'Etude Grokking} \hfill \textit{Juin--Septembre 2025} \\
    \textit{Remote} \hfill \textit{Python, PyTorch, TensorFlow, Math\'ematiques, R\'eseaux Neurones}
    \vspace{-2pt}
        \begin{itemize}[leftmargin=0.4cm, nosep, label=\tiny$\bullet$]
            \item \small{\'Etude g\'en\'eralisation tardive (Grokking) dans r\'eseaux neurones profonds, reproduction expériences SOTA.}
            \item \small{Analyse math\'ematique convergence, optimiseurs, dynamique apprentissage, interpr\'etabilit\'e mod\`eles.}
            \item \small{Imp\'lmentation PyTorch/TensorFlow, validation rigoureuse, contribution recherche publii\'ee.}
            \item \small{R\'edaction rapports techniques anglais, pr\'esentations \'equipe multidisciplinaire.}
        \end{itemize}
\end{itemize}

\section{Distinctions \& Leadership}
\begin{itemize}[leftmargin=*, label=\tiny$\bullet$, nosep]
    \item \small{\textbf{1\textsuperscript{er} prix IndabaX Africa 2025} (IA Sant\'e, hackathon panafricain)}
    \item \small{\textbf{1\textsuperscript{er} prix Mountain Hub R\'egional 2024} (Innovation technologique, IoT)}
    \item \small{\textbf{1\textsuperscript{er} prix CITS National 2024} (Smart City, ML pr\'edictif) -- vs 75 \'equipes}
    \item \small{\textbf{Head of Projects Cell \& Lead AI Cell ENSEA} -- +50\% membres actifs, collaborations Google DeepMind}
    \item \small{\textbf{Open Source :} Laravel API Generator (+30 \'etoiles GitHub, Packagist verified)}
    \item \small{\textbf{Certifications :} Supervised ML Stanford/DeepLearning.AI (Oct 2024), Python for Data Science IBM (Jan 2024)}
\end{itemize}
\end{document}
